\subsection{Problem definition and suburb selection}

An agency needs a tool estimating sale price from observable characteristics:
supervised regression, predicting $y \in \mathbb{R}^{+}$ from features $\mathbf{x}$.
The framing is hedonic \cite{rosen1974hedonic,malpezzi2003hedonic}. Accuracy alone is
insufficient, because the agency also needs to know when to distrust a prediction,
which motivates Section~4.

Three suburbs were chosen to span deliberately different markets, so location would be
a first-order predictor. They differ in level, spread and stock (Table~\ref{tab:suburbs}), i.e. Blacktown house-dominant and affordable, Chatswood
unit-dominant, Mosman mixed and premium. Suburb and type are entangled by construction,
so both are retained as inputs.

\begin{table}[H]
\centering\small
\begin{tabularx}{\textwidth}{@{}l XXX@{}}
\toprule
 & \textbf{Blacktown} & \textbf{Chatswood} & \textbf{Mosman} \\
\midrule
Retained sales $n$   & 79 & 62 & 47 \\
Median price         & \$950{,}000 & \$1{,}417{,}500 & \$1{,}760{,}000 \\
Mean price           & \$901{,}570 & \$1{,}716{,}668 & \$2{,}737{,}883 \\
Std.\ deviation      & \$334{,}345 & \$970{,}869 & \$2{,}116{,}702 \\
Range                & \$0.41--1.88M & \$0.70--4.80M & \$0.70--8.70M \\
Houses / Apartments  & 47 / 25 & 13 / 48 & 11 / 28 \\
Other stock          & 6 Twn, 1 Vil & 1 Twn & 4 Twn, 2 Ret, 2 Semi \\
\bottomrule
\end{tabularx}
\caption{Suburb contrast, computed in the notebook from the retained sample
($n=188$). Twn = townhouse, Ret = retirement living, Semi = semi-detached.}
\label{tab:suburbs}
\end{table}

\subsection{Collection, quality and limitations}

Records came from public \texttt{domain.com.au} sold listings \cite{domain2026sold},
with context from Domain profiles \cite{domain2026suburb} and ABS SEIFA
\cite{abs2021seifa}. realestate.com.au was excluded, as it denied programmatic access and
REA Group has an enforcement record \cite{om2024rea}. Every row carries
\texttt{source\_url} and \texttt{retrieved\_at} for verification purpose.

The original collection comprises of 239 properties, with missing values present in different columns.
Therefore, we aggressively discard incomplete rows, such that, there are more than 100 properties left with at least 30 properties in each suburb left.
Having removed the incomplete rows, \texttt{area\_sqm} is the only single source of missingness, so it
alone was exempted. Its remaining 64 gaps (34\%) are imputed inside the modelling pipeline.
Fortunately, three of those six columns are not model inputs, and it cost 27 Mosman sales, contributing to the sparse Mosman tail in the price distribution. Cross-validating
all three models on both the full and filtered sets confirms it is nonetheless the better choice, rather than keeping the missing values and conduct imputation in the modelling pipeline.

Noticeably, there are three apparent limitations. First, sales span only March--August 2026 with uneven coverage, so
the temporal term cannot be extrapolated. Second, \texttt{area\_sqm} remains $34\%$ missing, so that column depends on imputation.
Third, the SEIFA and \texttt{suburb\_*} columns resolve to three distinct rows, restating the suburb label rather than adding property-level information.

We also identify a potential source of bias is selection on price disclosure. Collection filtered to
listings showing a price, so withheld-price sales are absent by construction, and withholding is not random. The sample therefore under-represents the top of each distribution, most acutely in Mosman.
Recovering those sales requires the NSW Valuer General's statutory register \cite{vgnsw2026psi}, which records every transfer regardless of advertised disclosure.
